%% ==========================================================================
%% SECTION 1: INTRODUCTION (COMPOSED)
%% Skeleton: Version D (narrative arc, won 3/3 narrative reviewers)
%% Enrichment: Version B (Paxos framing, won 4 first-place votes from
%%   systems and ML reviewers)
%% Composition principle: D's linear arc with B's Paxos insight woven in
%% after the core argument is established.
%% ==========================================================================
\section{Introduction}\label{sec:intro}

%% --- CONTEXT: What world does the reader need to understand? ---

Systems that incorporate language model components face a verification
problem. The components generate text by sampling from learned probability
distributions, and some fraction of that text is fabricated: fluent,
confident, and wrong. We prefer the term \emph{fabrication} to emphasize
that fluent confabulation is the system's default behavior under coherence
optimization, not a malfunction to be patched. A language model optimized
for coherence produces coherent output. When the training data provides
grounding, that output tends to track truth. When it does not, the same
optimization produces fluent, confident fabrication.

Verification is not free. Checking a citation against a database costs API
calls and latency. Having a domain expert review a medical summary costs
time and salary. Running a second model as a judge costs compute. Every
system that deploys a language model component makes, explicitly or
implicitly, a budget allocation decision: how much of the output to verify,
which outputs to prioritize, and what signals to use for triage.

%% --- PROBLEM: What goes wrong? ---

\fakepara{The text-only interface cannot self-verify.}
The standard interface between a language model and its consumers exports
only text. A supervisor operating through this interface has exactly one
source of information: what the model chose to say---its \emph{testimony}.
But testimony is unreliable. Across four model architectures (OLMo-3 7B,
Llama-3.1 8B, Qwen3 4B, Mistral 7B), we find that every model reports
\emph{higher} confidence on fabrications than on known facts. Self-report
AUC ranges 0.28--0.36---below random, because the relationship is
systematically inverted. The models do not merely fail to report their
epistemic state; they systematically misreport it. This inversion is
universal: it holds across all four architectures, ruling out
training-procedure artifacts.

The problem is not limited to self-report. Stacking text-only supervisors
does not escape the constraint---later judges see no finer information than
earlier ones. A bounded supervisor receiving only text must spend its
verification budget reconstructing, from the output alone, what the model's
own computation already knew. The budget is consumed by reconstruction
rather than concentrated on the cases where verification is most needed.

%% --- CHALLENGES: Why is the problem hard? ---

\fakepara{Why existing approaches fall short.}
Xu et al.~\cite{xu2024hallucination} prove a related result: hallucination
is inevitable for LLMs over the class of computable functions. Our result is
orthogonal. They prove a computational limit on what the model can
\emph{produce}; we prove an observational limit on what a supervisor can
\emph{verify}. The model may know it is fabricating. The supervisor cannot
tell.

%% --- KEY INSIGHT: What changes the game? ---

\fakepara{Testimony versus telemetry.}
Language models carry internal signals---entropy traces, attention geometry,
token-level log-probabilities---that distinguish grounded generation from
fabrication. These signals are byproducts of the computation itself:
\emph{telemetry} (measurements of the computation that produced the output)
that the model cannot separately control under standard training. But the
standard text-only interface exports only the testimony and discards the
telemetry. Behavioral signals (response length, hedging frequency) can be
adapted by training without altering what the model knows, while
computational signals (per-token entropy, attention geometry) cannot be
flattened without changing the generation process itself.

%% --- PAXOS ENRICHMENT: Systems-level structural insight ---

This pattern maps to a well-known structure in distributed
systems. In Paxos consensus~\cite{paxos}, proposers generate candidate
values, acceptors filter proposals into decisions, and learners observe
only the decided values. A transformer's internal computation---feature
activations, competing attention heads, probability mass spread across the
vocabulary---plays the proposer role. Post-training alignment (RLHF, HHA)
plays the acceptor: it biases selection toward tokens that satisfy
helpfulness and coherence objectives. External observers are learners; they
see only the linearized token stream. The self-report inversion follows
directly: the acceptor optimizes for helpfulness, so when asked ``how
confident are you?'' it emits a confident answer, because confident answers
score well---regardless of how uncertain the proposal distribution was. In
standard Paxos, the acceptance criterion is fixed by protocol. In a
language model, the acceptor redefines its effective weights each round:
each output token changes the context, changing attention patterns, changing
the acceptance criterion for the next token. In Paxos this would be a
liveness bug; in LLMs it is the mechanism that enables both capability and
unreliability. Adding more learners---stacking supervisors---cannot recover
proposal history that the interface boundary already discarded.

%% --- SOLUTION: What did we build/prove? ---

\fakepara{Three results.}
This paper establishes three results that together characterize the cost
surface of epistemic verification.

First, we prove an \emph{impossibility}: epistemic honesty is unverifiable
for systems operating under text-only observation with bounded supervision
(\autoref{sec:formal}). The impossibility decomposes into two theorems. A
predictor-centric policy conditioning only on the query cannot satisfy
epistemic honesty across ambiguous world states
(\autoref{thm:representational}). Even with retrieval augmentation, a
bounded supervisor cannot provide the training signal needed to learn
honest behavior because plausible fabrications are observationally
indistinguishable from truthful responses (\autoref{thm:learnability}).
This impossibility functions as a budget constraint: below this investment
line, text-channel verification cannot help.

Second, we construct a \emph{tensor interface} that escapes the
impossibility class by exporting computational signals alongside text
(\autoref{sec:escape}). The interface requires no architectural changes to
current transformers. It exports entropy traces and attention
summaries---telemetry that the model cannot independently control. Under
bounded verification, tensor-guided judges at 10\% budget outperform
text-guided judges at 30\% budget (82.1\% vs.\ 78.5\% accuracy), not
because the tensor signal is strong in isolation, but because it enables
efficient triage: concentrating verification on the cases most likely to
contain fabrications.

Third, we provide an \emph{honest decomposition} of where the
discrimination comes from (\autoref{sec:eval}). The text channel is
stronger than expected: response length alone achieves AUC 0.88--0.96
across models---a free signal requiring no tensor interface. But it is
trivially gameable: a model trained to fabricate briefly defeats
length-based detection. After controlling for length, the tensor residual
adds AUC 0.60--0.70: a computational signal that cannot be independently
gamed without adversarial optimization targeting the signal itself. The gap
between the text-channel ceiling and the tensor-channel floor is the
impossibility theorem made empirical---it measures the boundary between
gameable and robust verification.

%% --- ACHIEVEMENTS: What does the reader gain? ---

\fakepara{A cost map, not a detector.}
The contribution of this paper is not a high AUC number. It is the map.
Systems engineers building with language model components need to know what
each level of verification investment buys, and where the investment
transitions from text-channel (cheap, gameable) to tensor-channel (more
expensive, robust). This paper provides that cost surface.

This paper makes the following contributions:

\begin{itemize}
\item \textbf{An impossibility result as budget constraint.}
  \autoref{sec:formal} proves that
  epistemic honesty is unverifiable for predictive token generation
  systems operating under text-only observation with bounded
  supervision. The impossibility is structural: it holds regardless
  of model scale, training procedure, or prompt engineering, and
  cannot be escaped by stacking supervisors. It establishes the
  floor below which text-channel verification investment cannot improve
  assurance.

\item \textbf{A tensor interface as the next investment tier.}
  \autoref{sec:escape} demonstrates a minimal interface extension that exports
  epistemic telemetry alongside text with no architectural changes to
  current transformers. The tensor interface increases the observation
  budget available to external verifiers by exposing computational signals
  the text-only interface discards.

\item \textbf{An empirical cost surface.}
  \autoref{sec:eval} decomposes verification effectiveness into
  text-channel and tensor-channel components, characterizes the
  cost-robustness tradeoff, and identifies the format-dependent
  variation that determines optimal judge strategy per domain.
\end{itemize}
